\medskip

On désigne par $f$ la fonction définie sur l'intervalle $[0~;~\pi]$ par

\[f(x) = \e^x \sin (x).\]

On note $\mathcal{C}_f$ la courbe représentative de $f$ dans un repère.

\bigskip

\textbf{PARTIE A}

\medskip

\begin{enumerate}
\item
	\begin{enumerate}
		\item Démontrer que pour tout réel $x$ de l'intervalle $[0~;~\pi]$,\:

\[f'(x) = \e^x[\sin(x) + \cos(x)].\]

		\item Justifier que la fonction $f$ est strictement croissante sur l'intervalle 
$\left[0~;~\dfrac{\pi}{2}\right]$
	\end{enumerate}
\item
	\begin{enumerate}
		\item Déterminer une équation de la tangente $T$ à la courbe $\mathcal{C}_f$ au point d'abscisse~$0$.
		\item Démontrer que la fonction $f$ est convexe sur l'intervalle $\left[0~;~\dfrac{\pi}{2}\right]$.
		\item En déduire que pour tout réel $x$ de l'intervalle $\left[0~;~\dfrac{\pi}{2}\right]$,\: $\e^x \sin (x) \geqslant x$.
	\end{enumerate}
\item Justifier que le point d'abscisse $\dfrac{\pi}{2}$ de la courbe représentative de la fonction $f$ est un point d'inflexion.
\end{enumerate}

\bigskip

\textbf{PARTIE B}

\medskip


On note

\[I = \displaystyle\int_0^{\frac{\pi}{2}} \e^x \sin (x)\, \text{d}x\quad  \text{et}\quad J = \displaystyle\int_0^{\frac{\pi}{2}} \e^x \cos (x)\, \text{d}x.\]

\smallskip

\begin{enumerate}
\item En intégrant par parties l'intégrale $I$ de deux manières différentes, établir les deux
 relations suivantes :

\[I= 1+J \qquad \text{et}\qquad I= \e^{\frac{\pi}{2}} - J.\]

\item En déduire que $I = \dfrac{1 + \e^{\frac{\pi}{2}}}{2}$.
\item On note $g$ la fonction définie sur $\R$ par $g(x) = x$.

Les courbes représentatives des fonctions $f$ et $g$ sont tracées dans le repère orthogonal ci-dessous sur l'intervalle $[0~;~\pi]$.

Calculer la valeur exacte de l'aire du domaine hachuré situé entre les courbes $\mathcal{C}_f$ et $\mathcal{C}_g$ et les droites d'équation $x = 0$ et $x = \dfrac{\pi}{2}$.

\begin{center}
\psset{xunit=1.8cm,yunit=0.9cm,arrowsize=2pt 3}
\begin{pspicture}(-1,-1)(4,9)
\psaxes[linewidth=1.25pt,Dy=2,labelFontSize=\scriptstyle]{->}(0,0)(0,0)(4,9)
\psline[linestyle=dotted,linewidth=2pt](1.5708,0)(1.5708,1.5708)
\psline[linewidth=1.25pt](1.5708,1.5708)(1.5708,4.81)
\psplot[plotpoints=1500,linewidth=1.25pt,linecolor=blue]{0}{3.14159}{2.71828 x exp x 180 mul 3.14159 div sin mul}
\psplot[plotpoints=1500,linewidth=1.25pt]{0}{3.14159}{x}
\uput[ur](2.7,6.6){\blue $\mathcal{C}_f$}\uput[dr](2.54,2.6){$\mathcal{C}_g$}
\pscustom[fillstyle=vlines]{\psplot[plotpoints=1500,linewidth=1.25pt,linecolor=blue]{0}{1.5708}{2.71828 x exp x 180 mul 3.14159 div sin mul}\psline[linewidth=1.25pt](1.5708,4.81)(1.5708,1.5708)(0,0)}
\end{pspicture}
\end{center}
\end{enumerate}
